\chapter{Peripherally Inserted Central Catheter}

\section*{Overview}
Peripherally inserted central catheter (PICC) line placement inserts a long catheter through a vein in the arm to reach the large central veins near the heart for medication, nutrition, or blood draws. The exact approach, setting, and preparation depend on the indication, anatomy, and the patient's health.

\section*{Alternative Names}
PICC line; peripherally inserted central catheter; PICC placement

\section*{Principle}
A flexible catheter is inserted into a peripheral vein of the upper arm and advanced so its tip rests in the superior vena cava. The tip position provides reliable access for irritant or long-term therapies.

\section*{History}
PICC lines were developed in the 1970s to provide central venous access from the arm, reducing the risks of chest placement and repeated needle sticks.

\section*{Why It Is Done}
Performed for patients needing several weeks of IV therapy such as antibiotics, chemotherapy, or total parenteral nutrition, and when peripheral access is difficult.

\section*{Is This a Primary or Secondary Test or Procedure}
Primary vascular access for intermediate-duration central venous therapy.

\section*{How It Is Performed}
Performed at the bedside or in a procedure room with ultrasound guidance. The vein is accessed in the upper arm, the catheter is threaded to the cavoatrial junction, position is confirmed (often by ECG guidance or radiograph), and the line is secured and dressed.

\section*{Preparation}
Review of bleeding risk and arm history (mastectomy, pacemaker, prior surgery), and ultrasound assessment of the veins.

\section*{Contraindications and Precautions}
Local infection, severe coagulopathy, and unsuitable arm veins are relative contraindications. Precaution: use the non-dominant arm and avoid the side of a prior mastectomy.

\section*{Risks}
Bleeding, infection, thrombosis, arrhythmia from tip migration, catheter malposition or breakage, and nerve injury.

\section*{Aftercare and Recovery}
The line is flushed, and the site is cared for with sterile dressing changes. Weekly flushing maintains patency when not in use.

\section*{Outcome and Follow-up}
A confirmatory radiograph shows the tip at the cavoatrial junction. Clinical use and site checks confirm function.

\section*{Expected Results}
A patent, well-positioned PICC line with no infection, thrombosis, or mechanical complication.

\section*{Follow-up}
Routine site inspection and flushing; the line is removed when therapy is complete.

\section*{When to Seek Medical Care}
Follow the procedure team's specific instructions. Seek urgent medical assessment for severe or worsening pain, uncontrolled bleeding, fever or signs of infection, trouble breathing, chest pain, fainting, new weakness or confusion, inability to urinate, or any rapidly worsening symptom. The appropriate warning signs depend on the procedure and the patient's medical condition.

\section*{References}
\begin{enumerate}
  \item MedlinePlus. Peripherally inserted central catheter (PICC) line. U.S. National Library of Medicine; 2025.
  \item Current Medical Diagnosis and Treatment. McGraw-Hill; 2025.
\end{enumerate}

\clearpage
